\documentclass{article}
\usepackage[english]{babel}
\usepackage{amsmath,graphicx}
\usepackage{geometry}
\graphicspath{{notes/assets/}{assets/}}

\begin{document}
\small

I would answer questions 1 and 3 together, since they are closely
related.

The VAE in our model is mainly for representation learning and compression. It gives a lower-dimensional latent space where manipulation is easier, but
direct VAE sampling is usually not accurate enough for single-cell simulation.
In principle, a standard VAE samples $z \sim \mathcal{N}(0,I)$ and then
$\hat{x} \sim p_\theta(x\mid z)$. This assumes that the decoder $p_\theta$ has learned a
good map from a Gaussian prior to the data distribution. In practice,
the encoded single-cell latents are much more complicated and sampling directly from
$\mathcal{N}(0,I)$ can therefore land in regions that do not match the
empirical latent distribution well. The diffusion model is added to learn this
latent distribution more accurately. This is analogous to latent diffusion in
image generation: the VAE learns a compressed representation, and the
diffusion model learns the distribution of that representation.

This also explains how our goal differs from scVI. scVI is mainly designed for representation learning and batch correction. It
can produce posterior samples, but those are not from the standard generation procedure discussed above. They are closer to
reconstruction than to genuine simulation. If $x$ is an
observed cell, the provided interface \texttt{posterior\_predictive\_sample()} samples
from $p_\theta(\hat{x}\mid x)$, that is, the observed cell is first encoded and
then decoded. This can give very realistic-looking samples, but each synthetic
cell is still anchored to a real input cell. It cannot 
generate new datasets of arbitrary size without providing matching real
cells.

The more relevant comparison is sampling from the prior, which means drawing
$z \sim \mathcal{N}(0,I)$ and then decoding it with the scVI decoder. They did not provide any interface for this, but I managed to do it by calling some internal functions. It truns out that the sample quality is usually much worse, where the UMAP below illustrates this
point. This suggests that VAE alone is not
accurate enough as a sampler for the learned latent space.

On the other hand, with a latent diffusion model, we can
generate cells by sampling new latents and decoding them, but with much
better sample quality than direct sampling from the VAE. The main advantage is that we
aim for high-quality genuine simulation while also keeping latent-space
controls for signals such as batch effect strength. The disadvantage is that we need to train and sample from an additional diffusion model, and we have to working on log-normalized space instead of raw counts since the library size variations is very hard to learn with limited data.

In short, our model is not a replacement for scVI. scVI focuses on representation learning and batch correction, whereas our model is more focused on high-quality simulation with controllable latent factors, which is a different but complementary goal.

\begin{center}
    \includegraphics[width=0.8\textwidth]{scvi_umap.png}
\end{center}

For question 2, I would interpret the latent blocks as supervised subspaces. For example, we may
write $z=(z_{\mathrm{cell}}, z_{\mathrm{batch}}, z_{\mathrm{residual}})$.
During VAE training, classification heads encourage cell-type information to
concentrate in $z_{\mathrm{cell}}$ and batch information to concentrate in
$z_{\mathrm{batch}}$. This gives us handles for intervention: for example, we
can shift only the batch block to introduce a batch effect while keeping the
cell-type block fixed.

The number of blocks is determined by the factors we want to control or
protect. In the current project, the natural choices are cell type and batch labels. The dimension of each block is a hyperparameter:
large enough to encode the target factor, but not so large that it occupies most of the latent space. We choose it empirically by checking disentanglement: the
intended label should be predictable from its own block, but much less
predictable from the other blocks, while reconstruction and simulation quality
remain stable. Our current choice is 16 dimensions over a total of 128 dimensions for both cell type and batch labels, which seems to work well.
Best,

Langtian

\end{document}
